% =====================================================================
%  GS-Net group-meeting results: tables + figures (validated/positive only).
%  Compile with pdflatex (needs booktabs, multirow, graphicx).
%  Figures: run `python make_figs.py` first to produce the PDFs, then this
%  file \includegraphics them (conference-style matplotlib).
% =====================================================================
\documentclass[11pt]{article}
\usepackage[margin=1in]{geometry}
\usepackage{booktabs}
\usepackage{multirow}
\usepackage{amsmath}
\usepackage{graphicx}
\usepackage{subcaption}

\begin{document}

% ---------------------------------------------------------------------
% TABLE 1 : main results
% ---------------------------------------------------------------------
\begin{table}[t]
\centering
\caption{Main results. GS-Net used as plug-and-play 3DGS initialization.
CARLA is multi-seed (degenerate scene~310 excluded); Waymo wide20 is the
sparse wide-baseline regime (real data), standard densify-on.}
\label{tab:main}
\begin{tabular}{llcccl}
\toprule
Dataset & Setting & Baseline & GS-Net & $\Delta$PSNR & Note \\
\midrule
CARLA & SSE (same-sensor)            & --    & --    & $\mathbf{+1.69\pm0.38}$ & reproduces paper ($+2.08$) \\
CARLA & CSE (cross-sensor, d.\ 2000) & 18.00 & 19.89 & $\mathbf{+1.89\pm0.10}$ & cross-sensor reuse \\
Waymo & SSE wide-baseline (real)     & 18.43 & 18.78 & $\mathbf{+0.35}$        & SSIM/LPIPS also better \\
\bottomrule
\end{tabular}
\end{table}

% ---------------------------------------------------------------------
% TABLE 2 : Waymo wide20 per-scene (densify-on)
% ---------------------------------------------------------------------
\begin{table}[t]
\centering
\caption{Waymo wide-baseline (wide20) SSE, per scene, standard densify-on.
GS-Net improves PSNR, SSIM and LPIPS.}
\label{tab:waymo}
\begin{tabular}{llccc}
\toprule
Scene & Method & PSNR & SSIM & LPIPS \\
\midrule
\multirow{2}{*}{Scene A}
 & baseline & 19.64 & 0.704 & 0.410 \\
 & GS-Net   & \textbf{20.34} & \textbf{0.723} & \textbf{0.384} \\
\midrule
\multirow{2}{*}{Scene B}
 & baseline & 17.23 & 0.639 & 0.482 \\
 & GS-Net   & 17.22 & 0.634 & 0.481 \\
\midrule
\multirow{2}{*}{\textbf{Avg}}
 & baseline & 18.43 & 0.671 & 0.446 \\
 & GS-Net   & \textbf{18.78} & \textbf{0.678} & \textbf{0.432} \\
\bottomrule
\end{tabular}
\end{table}

% ---------------------------------------------------------------------
% TABLE 3 : Waymo init spectrum (sparse wide-baseline) -- GS-Net is useful
% ---------------------------------------------------------------------
\begin{table}[t]
\centering
\caption{Waymo sparse wide-baseline init-spectrum (SSE PSNR): SfM-init /
MVS-init (densification ceiling) / GS-Net-init. The densification ceiling is
large ($+2.40$) and GS-Net captures part of it ($+0.64$ over SfM), i.e.\ GS-Net
is useful where real coverage holes exist.}
\label{tab:spectrum}
\begin{tabular}{lcccc}
\toprule
Init & SfM & MVS (ceiling) & GS-Net & GS-Net $-$ SfM \\
\midrule
SSE PSNR & 18.03 & 20.44 & 18.68 & $\mathbf{+0.64}$ \\
\bottomrule
\end{tabular}
\end{table}

% ---------------------------------------------------------------------
% TABLE 4 : encoder ablation
% ---------------------------------------------------------------------
\begin{table}[t]
\centering
\caption{Encoder ablation (CARLA, final loss recipe, 4 sequences). The top-4
encoders are flat within noise ($0.32$\,dB): the gain comes from the loss
design, not encoder complexity $\Rightarrow$ a lightweight encoder suffices.}
\label{tab:encoder}
\begin{tabular}{lcc}
\toprule
Encoder & PSNR & $\Delta$ vs baseline \\
\midrule
concat (light)   & \textbf{26.65} & $+2.09$ \\
geom             & 26.51 & $+1.95$ \\
mlp-only         & 26.40 & $+1.84$ \\
attention        & 26.33 & $+1.77$ \\
edgeconv         & 25.37 & $+0.81$ \\
\bottomrule
\end{tabular}
\end{table}

% ---------------------------------------------------------------------
% FIGURES  (run make_figs.py to produce the PDFs)
% ---------------------------------------------------------------------
\begin{figure}[t]
\centering
\begin{subfigure}{0.48\linewidth}
  \includegraphics[width=\linewidth]{fig_tsweep.pdf}
  \caption{Sensitivity to $T$ (CARLA). $T{=}5$ best; larger $T$ no gain.}
  \label{fig:tsweep}
\end{subfigure}\hfill
\begin{subfigure}{0.48\linewidth}
  \includegraphics[width=\linewidth]{fig_densify.pdf}
  \caption{CSE gain vs.\ densification budget; peaks at a moderate budget.}
  \label{fig:densify}
\end{subfigure}
\caption{Sensitivity studies. (\subref{fig:tsweep}) number of expansion heads;
(\subref{fig:densify}) cross-sensor gain is largest at a moderate densification
budget ($\approx2000$): GS-Net's view-consistent init reaches quality early
while the baseline needs heavy densification to catch up.}
\label{fig:sens}
\end{figure}

\begin{figure}[t]
\centering
\includegraphics[width=0.82\linewidth]{fig_pseudogt.pdf}
\caption{Robustness to pseudo-GT quality (Reviewer~C). GS-Net degrades
gracefully along both axes -- weaker supervision and sparser pseudo-GT --
with no collapse.}
\label{fig:pseudogt}
\end{figure}

\end{document}
